\chapter{Bridging the Intensional--Extensional Gap}
\label{ch:fallback}

The preceding chapters built a powerful \emph{intensional}
equivalence engine: OTerm normalization reduces programs to canonical
forms, the HIT structural prover applies path constructors by
structural decomposition, and bidirectional BFS searches for
multi-step rewrite paths using 16 registered axiom functions.
The Z3 per-fiber checker then provides a separate layer of
verification via \v{C}ech cohomology on the type fibration.

This chapter examines \emph{where this pipeline reaches its limits},
\emph{why} those limits arise from a fundamental gap in type theory,
and \emph{how CCTT provides a principled programme for extending the
pipeline} by enriching the path space, strengthening specification
identification, and improving the interaction between denotational
and cohomological layers.  The chapter is forward-looking: it
identifies the theoretical mechanisms that motivate each next
improvement to the implementation.


%% ===============================================================
\section{The Current Pipeline as a Typed Function Chain}
\label{sec:current-pipeline}
%% ===============================================================

The implementation orchestrates four strategy layers, each with
clear CCTT semantics.  We formalize the full pipeline as a chain
of typed functions with a soundness invariant.

\begin{definition}[Pipeline Type]
\label{def:pipeline-type}
Define the type of pipeline verdicts:
\[
  \mathsf{Verdict} \coloneqq \mathsf{EQ} \mid \mathsf{NEQ} \mid
  \mathsf{UNKNOWN}(c)
  \qquad c \in [0,1]
\]
and the pipeline as a composition of partial decision procedures:
\[
  \Pi \coloneqq \pi_4 \circ \pi_3 \circ \pi_2 \circ \pi_1 \circ \pi_0
  : (\mathsf{Src} \times \mathsf{Src}) \to \mathsf{Verdict}
\]
where each $\pi_k$ is a function
$\mathsf{Verdict}_? \to \mathsf{Verdict}_?$
that either refines the verdict or passes it through unchanged,
and $\mathsf{Verdict}_?$ extends $\mathsf{Verdict}$ with an
explicit $\bot$ (``no verdict yet'').
\end{definition}

The layers are:

\begin{enumerate}
  \item[\textbf{$\pi_0$}:] \textbf{Closed-term evaluation}
    (zero-argument functions).
    A function with no parameters has a singleton domain; its
    denotation IS its unique value.  Evaluating it computes
    a $\beta$-normal form --- not a sample.  Functions with only
    variadic parameters (\texttt{*args}) are evaluated at zero
    arguments as a \emph{witness}: trusted only for NEQ (a single
    disagreement suffices) but not for EQ (agreement on one input
    is insufficient).

  \item[\textbf{$\pi_1$}:] \textbf{Denotational OTerm equivalence}
    (the primary weapon).  Both functions are compiled to OTerms
    via \texttt{compile\_to\_oterm}, passed through the 7-phase
    normalizer (\S\ref{sec:normalizer-formal}), and compared.
    Four sub-strategies operate in sequence:
    \begin{enumerate}[label=(\alph*)]
      \item \emph{Exact canonical form match}: after normalization,
        the string representations are identical.  Five gating
        conditions must hold: (i) canonical form length $\geq 20$
        or constant value, (ii) no \texttt{OUnknown} nodes, (iii)
        all free variables are parameters, (iv) at least one
        parameter referenced, (v) operation-name sets agree.
      \item \emph{Flexible matching}: canonical forms differ only
        in OFix keys (hash-based loop body identifiers) but are
        otherwise structurally identical under
        \texttt{\_flexible\_match}.
      \item \emph{HIT path induction + BFS path search}
        (\S\ref{sec:hit-rewrite-system}): the elimination
        principle for $\mathsf{PyComp}$.  Structural decomposition
        with congruence rules at each OTerm constructor, cross-type
        path constructors (D1--D24), and bidirectional BFS with
        16 axiom functions.  Fiber-aware via duck-type context.
      \item \emph{Provable NEQ detection}
        (\S\ref{sec:neq-formal}): structural impossibilities
        (different effect signatures, incompatible literal values,
        non-commutative argument orders on fiber-restricted
        operations).
    \end{enumerate}

  \item[\textbf{$\pi_2$}:] \textbf{Z3 structural equality}: the
    Z3 presheaf sections (guard--term pairs) are compared
    syntactically, but only trusted when \emph{fully interpreted}
    (no uninterpreted function symbols), as checked by
    \texttt{\_uninterp\_depth}.

  \item[\textbf{$\pi_3$}:] \textbf{Z3 per-fiber checking +
    \v{C}ech $H^1$}: duck-type inference determines the site
    category, the system checks equivalence on each fiber with
    bounded timeout, and \v{C}ech cohomology assembles the local
    judgments into a global verdict.  $H^1 = 0$ implies EQ;
    $H^1 \neq 0$ implies NEQ.
\end{enumerate}

\begin{theorem}[Pipeline Soundness --- Informal]
\label{thm:pipeline-soundness-informal}
The pipeline $\Pi$ satisfies: if $\Pi(f,g) = \mathsf{EQ}$
then $\forall x.\; f(x) = g(x)$; if $\Pi(f,g) = \mathsf{NEQ}$
then $\exists x.\; f(x) \neq g(x)$.  A full proof appears in
\S\ref{sec:fallback-soundness}.
\end{theorem}

The key architectural invariant: each $\pi_k$ is
\emph{short-circuiting}: if a previous layer has already produced
$\mathsf{EQ}$ or $\mathsf{NEQ}$, the remaining layers are skipped.
The ordering is calibrated so cheaper, more precise layers run first.

This pipeline currently proves 7/50 equivalences and 48/50
non-equivalences on the hard benchmark suite.


%% ===============================================================
\section{The 7-Phase Normalizer}
\label{sec:normalizer-formal}
%% ===============================================================

The normalizer is the engine behind sub-strategy (a) of $\pi_1$.
It applies seven phases in a fixed-point loop: the entire sequence
is iterated until the canonical form stabilizes (up to 8 iterations).

\begin{definition}[Normalizer]
\label{def:normalizer}
$\mathcal{N} \coloneqq (\varphi_7 \circ \varphi_6 \circ \varphi_5
\circ \varphi_4 \circ \varphi_3 \circ \varphi_2 \circ \varphi_1)^{\leq 8}$,
where $t^{\leq k}$ means ``apply $t$ until fixpoint, at most $k$
iterations.''
\end{definition}

The seven phases are:

\begin{center}
\begin{tabular}{clp{6.5cm}}
\hline
\textbf{Phase} & \textbf{Name} & \textbf{Key Rules} \\
\hline
$\varphi_1$ & $\beta$-reduction
  & Constant folding, $\mathsf{case}(\top, t, f) \to t$,
    $\mathsf{case}(c, x, x) \to x$, $\fold(\mathit{op}, i, []) \to i$,
    identity-map elimination \\
$\varphi_2$ & Ring/lattice
  & Additive/multiplicative identities, double negation, sorted
    idempotence, comparison negation \\
$\varphi_3$ & Fold canon.
  & $\mathsf{sum}(x) \to \fold(\mathsf{iadd}, 0, x)$,
    $\mathsf{any}(x) \to \fold(\lor, \bot, x)$,
    $\mathsf{join} \to \fold(\mathsf{str\_concat})$ \\
$\varphi_4$ & HOF fusion
  & $\mathsf{sorted}(\mathsf{list}(x)) \to \mathsf{sorted}(x)$,
    $\mathsf{reversed}(\mathsf{sorted}(x)) \to \mathsf{sorted\_rev}(x)$,
    $\mathsf{abs}(x) \to \mathsf{case}(\mathsf{lt}(x,0), -x, x)$,
    $\mathsf{max}(a,b) \to \mathsf{case}(\mathsf{gte}(a,b), a, b)$ \\
$\varphi_5$ & Symbol unify
  & D19: $\mathsf{heapq.nsmallest} \to \mathsf{sorted\_slice}$;
    D13: $\mathsf{Counter} \to \mathsf{counter}$;
    D14: $\mathsf{OrderedDict} \to \mathsf{dict}$;
    local variable method normalization \\
$\varphi_6$ & Quotient
  & $Q[\mathsf{perm}](\mathsf{sorted}(x)) \to \mathsf{sorted}(x)$,
    $Q[Q[\cdot]] \to Q[\cdot]$ (idempotence),
    $\mathsf{sorted}(Q[\mathsf{perm}](x)) \to \mathsf{sorted}(x)$ \\
$\varphi_7$ & Piecewise
  & Flatten nested case chains, infer complement guards via
    comparison algebra, sort pieces by
    $(\mathsf{value.canon()}, \mathsf{guard.canon()})$ \\
\hline
\end{tabular}
\end{center}

\subsection{Confluence and Termination}

\begin{proposition}[Termination]
\label{prop:normalizer-termination}
The normalizer terminates in at most 8 iterations.  Each phase is
a structural recursion on the OTerm tree (hence terminates on each
invocation), and the outer loop is bounded by the explicit iteration
cap.
\end{proposition}

\begin{proof}
Each phase traverses the OTerm by structural recursion (depth-first,
polynomial in term size).  The outer loop checks
$\mathsf{canon}(t_{\mathit{prev}}) = \mathsf{canon}(t_{\mathit{curr}})$
and exits early when the fixpoint is reached.  The hard bound of 8
iterations ensures termination even if no fixpoint is reached.
\end{proof}

\begin{proposition}[Weak Confluence]
\label{prop:normalizer-confluence}
Within each phase, the rewrite rules are \emph{non-overlapping}
at the root: each rule matches a distinct OTerm constructor or
operation name.  Consequently, each phase is confluent on its own
inputs.
\end{proposition}

The cross-phase interaction is more subtle: $\varphi_3$ may create
$\fold$ nodes that $\varphi_4$ then fuses, which $\varphi_6$ then
quotients.  Confluence of the composed system $\varphi_7 \circ
\cdots \circ \varphi_1$ is not proved in general but is empirically
observed: all benchmark pairs converge within 3 iterations.

\begin{proposition}[Soundness of Each Phase]
\label{prop:phase-soundness}
For each phase $\varphi_k$: if $\varphi_k(t) = t'$ then
$\den{t} = \den{t'}$ (extensional equality of denotations).
\end{proposition}

\begin{proof}[Proof sketch]
Phase by phase:
$\varphi_1$ applies standard $\beta$-reduction and constant folding
(sound by substitution);
$\varphi_2$ applies ring identities ($x + 0 = x$, $x \cdot 1 = x$,
etc.) that are valid in $\mathbb{Z}$, $\mathbb{R}$, and the
Boolean algebra;
$\varphi_3$ replaces builtins by their fold definitions (sound by
the semantics of \texttt{sum}, \texttt{any}, etc.);
$\varphi_4$ applies HOF fusion laws ($\mathsf{sorted} \circ
\mathsf{list} = \mathsf{sorted}$, etc.) that hold in Python;
$\varphi_5$ unifies module-qualified names to canonical forms (same
function, different import path);
$\varphi_6$ exploits the quotient type: $Q[\mathsf{perm}]$ marks
nondeterministic ordering, and $\mathsf{sorted}$ eliminates it;
$\varphi_7$ reorders case-guards, valid because piecewise functions
over a partition are determined by their values on each piece.
\end{proof}


%% ===============================================================
\section{The HIT Prover as a Term-Rewriting System}
\label{sec:hit-rewrite-system}
%% ===============================================================

The HIT prover (\texttt{hit\_path\_equiv}) implements the
elimination principle for $\mathsf{PyComp}$
(Theorem~\ref{thm:pycomp-elim}).  We now formalize it as a
\emph{conditional term-rewriting system} (CTRS) and establish its
rewrite-theoretic properties.

\subsection{Congruence Rules}

At each OTerm constructor, the prover applies a \emph{congruence rule}:
decompose the top-level constructor and recurse into children.
These rules form the structural core of the prover:

\[
\begin{array}{rl}
\textsc{Refl}: & \displaystyle
  \frac{a.\mathsf{canon}() = b.\mathsf{canon}()}
       {a =_{\mathsf{HIT}} b} \\[12pt]
\textsc{OOp-Cong}: & \displaystyle
  \frac{a_1 =_{\mathsf{HIT}} b_1 \quad \cdots \quad a_n =_{\mathsf{HIT}} b_n}
       {\mathsf{OOp}(f, a_1, \ldots, a_n) =_{\mathsf{HIT}} \mathsf{OOp}(f, b_1, \ldots, b_n)} \\[12pt]
\textsc{OCase-Cong}: & \displaystyle
  \frac{g =_{\mathsf{HIT}} g' \quad t =_{\mathsf{HIT}} t' \quad f =_{\mathsf{HIT}} f'}
       {\mathsf{case}(g, t, f) =_{\mathsf{HIT}} \mathsf{case}(g', t', f')} \\[12pt]
\textsc{OCase-Flip}: & \displaystyle
  \frac{g' = \lnot g \quad t =_{\mathsf{HIT}} f' \quad f =_{\mathsf{HIT}} t'}
       {\mathsf{case}(g, t, f) =_{\mathsf{HIT}} \mathsf{case}(g', t', f')} \\[12pt]
\textsc{OFold-Cong}: & \displaystyle
  \frac{\mathit{op} = \mathit{op}' \quad i =_{\mathsf{HIT}} i' \quad c =_{\mathsf{HIT}} c'}
       {\fold(\mathit{op}, i, c) =_{\mathsf{HIT}} \fold(\mathit{op}', i', c')} \\[12pt]
\textsc{OLam-Cong}: & \displaystyle
  \frac{|\vec{x}| = |\vec{y}| \quad \mathit{body}_a[\vec{x}/\vec{y}] =_{\mathsf{HIT}} \mathit{body}_b}
       {\lambda\vec{x}.\mathit{body}_a =_{\mathsf{HIT}} \lambda\vec{y}.\mathit{body}_b} \\[12pt]
\textsc{OMap-Cong}: & \displaystyle
  \frac{t =_{\mathsf{HIT}} t' \quad c =_{\mathsf{HIT}} c' \quad p =_{\mathsf{HIT}} p'}
       {\mathsf{map}(t, c, p) =_{\mathsf{HIT}} \mathsf{map}(t', c', p')} \\[12pt]
\textsc{Comm}: & \displaystyle
  \frac{\mathit{op} \in \mathsf{Comm}_\tau \quad a =_{\mathsf{HIT}} b' \quad b =_{\mathsf{HIT}} a'}
       {\mathsf{OOp}(\mathit{op}, a, b) =_{\mathsf{HIT}} \mathsf{OOp}(\mathit{op}, a', b')}
\end{array}
\]

The $\textsc{Comm}$ rule is \emph{fiber-restricted}: the set
$\mathsf{Comm}_\tau$ of commutative operations depends on the duck-type
fiber $\tau$.  On numeric fibers, $\mathsf{add} \in \mathsf{Comm}_\tau$;
on string/list fibers, it is not (concatenation is non-commutative).
This is sheaf descent (\S2.6) at the rewriting level.

\subsection{Cross-Type Path Constructors}

Beyond congruence, the prover applies cross-type rules that relate
OTerms with different top-level constructors:

\begin{itemize}
  \item \textbf{D1}: $\mathsf{OFix} \leftrightarrow \mathsf{OFold}$
    (recursion $\leftrightarrow$ iteration) --- via spec matching
    on the fix-point body.
  \item \textbf{D8}: Case flattening ---
    $\mathsf{case}(t_1, A, \mathsf{case}(t_2, B, C))
    \leftrightarrow \mathsf{case}(t_2, B, \mathsf{case}(t_1, A, C))$
    when guards are disjoint.
  \item \textbf{D20}: Spec identification (Yoneda) --- both sides
    reduce to the same $\mathsf{OAbstract}(\mathit{spec}, \vec{x})$.
  \item \textbf{D22}: $\mathsf{OCatch} \leftrightarrow \mathsf{OCase}$
    (try/except $\leftrightarrow$ conditional) --- matching body and
    default.
  \item \textbf{De Morgan}:
    $\lnot(a \land b) \leftrightarrow (\lnot a) \lor (\lnot b)$,
    applied recursively.
  \item \textbf{Guard subsumption}:
    if $G_{\mathit{sub}}$ is a disjunct of $G$, then in the
    else-branch of $G$, the inner
    $\mathsf{case}(G_{\mathit{sub}}, \ldots)$ collapses.
  \item \textbf{Comparison normalization}:
    $\mathsf{lt}(a,b) = \mathsf{gt}(b,a)$,
    $\lnot\mathsf{lt}(a,b) = \mathsf{gte}(a,b)$, etc.
\end{itemize}

The prover also applies single-step axiom rewrites from the
16 registered axiom functions at depths $< 8$, checking whether
any one-step rewrite closes the gap.

\subsection{Memoization and Termination}

The HIT prover uses a memoization table keyed by
$(\mathsf{canon}(a), \mathsf{canon}(b))$.  Before recursing,
it marks the key as ``in progress'' to detect cycles (recursive
OTerms produce cycles in the dependency graph).  A depth bound
of 20 ensures termination.

\begin{proposition}[HIT Prover Termination]
The HIT prover terminates on all inputs.  The depth bound and
memoization ensure each pair is visited at most once per depth
level, giving worst-case $O(|t|^2 \cdot D)$ where $|t|$ is
the OTerm size and $D = 20$ is the depth bound.
\end{proposition}


%% ===============================================================
\section{The D20 Spec Identification Engine}
\label{sec:d20-formal}
%% ===============================================================

The D20 axiom function (\texttt{\_axiom\_d20\_spec\_unify})
recognizes high-level specifications from OTerm structure and
replaces complex computations with abstract specification terms.
This is the Yoneda embedding: a program is characterized by
the specification it satisfies.

\subsection{Spec Identification as a Morphism}

\begin{definition}[Spec Identification Function]
\label{def:spec-id}
The spec identification function is a partial map:
\[
  \iota : \mathsf{OTerm} \rightharpoonup
  \mathsf{SpecName} \times \mathsf{OTerm}^*
\]
defined by structural pattern matching on the OTerm.
\end{definition}

The current implementation recognizes five specification families:

\begin{center}
\begin{tabular}{lp{6cm}l}
\hline
\textbf{Spec} & \textbf{Pattern} & \textbf{Result} \\
\hline
$\mathsf{factorial}$ &
  $\fold(\mathsf{mul}, 1, \mathsf{range}(\ldots, n))$ &
  $\mathsf{factorial}(n)$ \\
$\mathsf{range\_sum}$ &
  $\fold(\mathsf{add}, 0, \mathsf{range}(\ldots, n))$ &
  $\mathsf{range\_sum}(n)$ \\
$\mathsf{fib\_like}$ &
  $\mathsf{fix}[h](\ldots\mathsf{add}\ldots\mathsf{sub}\ldots)$ &
  $\mathsf{fib\_like}(\vec{v})$ \\
$\mathsf{sorted}$ &
  $\mathsf{sorted}(x)$ &
  $\mathsf{sorted}(x)$ \\
$\mathsf{binomial}$ &
  $\mathsf{math.comb}(n,k)$ &
  $\mathsf{binomial}(n,k)$ \\
\hline
\end{tabular}
\end{center}

\begin{proposition}[Spec Identification Soundness]
\label{prop:spec-id-sound}
If $\iota(t) = (S, \vec{x})$, then $t$ computes the specification
$S$ on inputs $\vec{x}$.  In particular, if $\iota(t_1) = (S, \vec{x})$
and $\iota(t_2) = (S, \vec{x})$ (same spec, same inputs), then
$\den{t_1} = \den{t_2}$ by Theorem~\ref{thm:det-spec}.
\end{proposition}

\begin{proof}
Each pattern in $\iota$ is a sufficient condition for the
corresponding specification.
$\fold(\mathsf{mul}, 1, \mathsf{range}(1, n{+}1))$ computes
$\prod_{i=1}^n i = n!$ by definition of fold over range.
$\fold(\mathsf{add}, 0, \mathsf{range}(\ldots))$ computes
$\sum_{i} i$ by the same argument.
$\mathsf{sorted}(x)$ is the identity on the sorted spec.
The Fibonacci pattern uses a heuristic (presence of \texttt{add}
and \texttt{sub} in the fix body), which is sound when the fix
body matches the Fibonacci recurrence but may produce false
negatives on non-standard formulations.
\end{proof}


\subsection{Formal Pattern Matching as Presheaf Morphism}

The spec identification function $\iota$ can be understood as a
natural transformation between presheaves.  Define:
\begin{itemize}
  \item $\mathcal{O}$ (the OTerm presheaf): assigns to each type
    fiber $\tau$ the set of OTerms whose free variables inhabit
    $\tau$.
  \item $\mathcal{S}$ (the spec presheaf): assigns to each $\tau$
    the set of specifications recognizable on $\tau$.
\end{itemize}

Then $\iota : \mathcal{O} \Rightarrow \mathcal{S}$ is a natural
transformation: for each restriction map $\rho_{\tau \to \tau'}$,
the diagram
\[
  \begin{tikzcd}
    \mathcal{O}(\tau) \ar[r, "\iota_\tau"] \ar[d, "\rho"'] &
    \mathcal{S}(\tau) \ar[d, "\rho"] \\
    \mathcal{O}(\tau') \ar[r, "\iota_{\tau'}"] &
    \mathcal{S}(\tau')
  \end{tikzcd}
\]
commutes: restricting an OTerm to a sub-fiber and then identifying
its spec gives the same result as identifying the spec and then
restricting.  This naturality ensures that spec identification
is compatible with the sheaf structure of the type fibration.


%% ===============================================================
\section{BFS Path Search and the 16 Axiom Functions}
\label{sec:bfs-axioms}
%% ===============================================================

When the HIT structural prover cannot close the gap directly,
the system falls back to bidirectional BFS on the rewrite
groupoid (\S\ref{sec:enriching-paths}).  The BFS uses 16
registered axiom functions, each of which generates all
one-step rewrites applicable to a given OTerm.

\subsection{The 16 Registered Axiom Functions}

The axiom functions are organized by the monograph's dichotomy
groups.  Each function has type
$\mathsf{OTerm} \times \mathsf{FiberCtx} \to
\mathsf{List}(\mathsf{OTerm} \times \mathsf{String})$,
returning all applicable rewrites with their axiom name:

\begin{center}
\begin{tabular}{clp{5.5cm}}
\hline
\textbf{\#} & \textbf{Name} & \textbf{Axiom} \\
\hline
1 & D1 fold/unfold & $\mathsf{fix} \leftrightarrow \fold$ \\
2 & D2 $\beta$ & $(\lambda x.b)(a) \to b[x := a]$ \\
3 & D4 map fusion & $\mathsf{map}(f, \mathsf{map}(g, c)) \to \mathsf{map}(f \circ g, c)$ \\
4 & D5 fold universal & Hash-named fold $\to$ canonical op \\
5 & D8 section merge & Guard complement collapse \\
6 & D13 histogram & $\mathsf{Counter}(x) \leftrightarrow \fold(\mathsf{count}, \ldots, x)$ \\
7 & D16 memo/DP & Recurrence normalization \\
8 & D17 associativity & Tree fold $\to$ linear fold \\
9 & D19 sort-scan & $\fold(\mathit{op}, \mathsf{sorted}(x)) \to \fold(\mathit{op}, x)$ when $\mathit{op}$ commutative \\
10 & D20 spec unify & $t \to \mathsf{OAbstract}(\mathit{spec}, \vec{x})$ \\
11 & D22 effect elim & $\mathsf{catch}(b, d) \leftrightarrow \mathsf{case}(g, b, d)$ \\
12 & D24 $\eta$ & $\lambda x.f(x) \to f$ \\
13 & Algebra & Commutativity, associativity, identity \\
14 & Quotient & $Q[\mathsf{perm}]$ absorption, idempotence \\
15 & Fold & Sum/prod expansion, range normalization \\
16 & Case & Constant guard, identical branches, $\lnot$-flip, subsumption \\
\hline
\end{tabular}
\end{center}

Each axiom function also operates under \emph{congruence closure}:
the function \texttt{\_congruence\_rewrites} applies every axiom
at every sub-term position, constructing the new OTerm by replacing
the rewritten sub-term and annotating the position (e.g.,
\texttt{D4\_map\_fusion@arg0}).

\subsection{Axiom Soundness}

\begin{theorem}[Axiom Soundness]
\label{thm:axiom-soundness}
Each of the 16 axiom functions preserves semantic equivalence:
if $(t', \mathit{ax}) \in \mathsf{axiom}_k(t, \mathit{ctx})$
then $\den{t} = \den{t'}$.
\end{theorem}

\begin{proof}
We verify soundness per axiom:
\begin{enumerate}
  \item \textbf{D1}: A fix-point with accumulation pattern computes
    the same function as the corresponding fold, by the universal
    property of folds (catamorphism).
  \item \textbf{D2}: $\beta$-reduction is sound by the substitution
    lemma.
  \item \textbf{D4}: Map fusion: $\mathsf{map}(f, \mathsf{map}(g, c))
    = \mathsf{map}(f \circ g, c)$ by the functor law.
  \item \textbf{D5}: Fold canonicalization replaces a hash-based fold
    key with a recognized canonical operation, preserving semantics.
  \item \textbf{D8}: Guard complement collapse:
    $\mathsf{case}(g, A, \mathsf{case}(\lnot g, B, C))
    = \mathsf{case}(g, A, C)$ because $\lnot g$ is false in the
    else-branch of $g$.
  \item \textbf{D13}: $\mathsf{Counter}(x)$ is defined as
    $\fold(\mathsf{count\_freq}, \mathsf{defaultdict}(0), x)$.
  \item \textbf{D16}: Recurrence normalization alpha-normalizes the
    fix body, preserving the function computed.
  \item \textbf{D17}: Tree fold = linear fold when the operation is
    associative (parenthesization doesn't matter).
  \item \textbf{D19}: Sorting is irrelevant for commutative folds:
    $\fold(\mathit{op}, i, \mathsf{sorted}(x))
    = \fold(\mathit{op}, i, x)$ when $\mathit{op}$ is commutative
    and associative.
  \item \textbf{D20}: Spec identification replaces a term with its
    abstract specification; soundness follows from
    Proposition~\ref{prop:spec-id-sound}.
  \item \textbf{D22}: When the exception guard is decidable,
    $\mathsf{catch}(b, d) = \mathsf{case}(\mathsf{can\_succeed}(b), b, d)$.
  \item \textbf{D24}: $\eta$-contraction:
    $\lambda x.f(x) = f$ when $x \notin \mathrm{FV}(f)$.
  \item \textbf{Algebra}: Standard algebraic identities
    (commutativity restricted to appropriate fibers,
    associativity, identity elements, involution).
  \item \textbf{Quotient}: $Q[\mathsf{perm}](\mathsf{sorted}(x))
    = \mathsf{sorted}(x)$ because sorting already fixes a
    canonical representative.
  \item \textbf{Fold}: $\mathsf{sum}(x) = \fold(\mathsf{add}, 0, x)$
    etc., by definition.
  \item \textbf{Case}: Constant guard elimination and branch
    identity are immediate; $\lnot$-flip swaps branches soundly.
\end{enumerate}
\end{proof}

\subsection{BFS Structure}

The bidirectional BFS maintains two frontiers:
$F_0 = \{\mathsf{nf}_f\}$ and $B_0 = \{\mathsf{nf}_g\}$.
At each depth $d$, both frontiers are expanded by applying all
16 axiom functions (with congruence closure) to all terms in the
current-depth layer.  Frontier pruning keeps each set below
$\mathsf{max\_frontier} = 200$ terms, scored by path length
and structural similarity to the target.

A path is found when a canonical form appears in both frontiers.
The maximum search depth is $\mathsf{max\_depth} = 4$,
giving a worst-case exploration of
$O(\mathsf{max\_depth} \cdot \mathsf{max\_frontier} \cdot
|\mathsf{Axioms}| \cdot |t|)$ term rewrites.


%% ===============================================================
\section{Provable NEQ Detection}
\label{sec:neq-formal}
%% ===============================================================

The function \texttt{\_detect\_provable\_neq} returns $\top$ only
when the two OTerms are \emph{provably} non-equivalent.  Each
detection rule is a certificate of non-equivalence.

\begin{definition}[NEQ Certificate]
\label{def:neq-cert}
A \emph{NEQ certificate} for OTerms $a, b$ is a structural
property $P(a, b)$ such that $P(a, b) \implies \exists x.\;
\den{a}(x) \neq \den{b}(x)$.
\end{definition}

The implementation checks the following certificates, organized
by category:

\paragraph{Effect-fiber certificates.}
If $a$ contains an effect operation (\texttt{effect\_io},
\texttt{effect\_write}, etc.) and $b$ does not (or vice versa),
they inhabit different fibers of the heap fibration:
$H^1 \neq 0$ on the effect fiber.

\paragraph{Literal certificates.}
$\mathsf{OLit}(v_1)$ vs.\ $\mathsf{OLit}(v_2)$ with $v_1 \neq v_2$:
constant functions returning different values.
$\mathsf{OLit}(\mathsf{None})$ vs.\ any non-$\mathsf{None}$ term:
the denotation $\bot$ differs from any productive computation.

\paragraph{Argument-order certificates.}
For non-commutative operations
$\sigma \in \{\mathsf{sub}, \mathsf{floordiv}, \mathsf{mod},
\mathsf{pow}, \mathsf{concat}, \mathsf{merge}, \ldots\}$:
if $\sigma(a_1, a_2)$ vs.\ $\sigma(a_2, a_1)$ with
$a_1 \neq a_2$, then the results differ.
For \texttt{add}/\texttt{iadd}, commutativity depends on the
duck-type fiber: on numeric fibers, it is commutative (no
certificate); on string/list fibers, it is non-commutative
(certificate issued).

\paragraph{Comparison-operator certificates.}
Same arguments, different comparison:
$\mathsf{eq}(a, b)$ vs.\ $\mathsf{is}(a, b)$ differ because
\texttt{==} and \texttt{is} have different semantics in Python
(value vs.\ identity).

\paragraph{Structural certificates.}
\begin{itemize}
  \item Different sequence lengths:
    $|\mathsf{OSeq}(a)| \neq |\mathsf{OSeq}(b)|$.
  \item Different fold operators:
    $\fold(\mathit{op}_1, \ldots) \neq \fold(\mathit{op}_2, \ldots)$
    when $\mathit{op}_1 \neq \mathit{op}_2$.
  \item Cross-type at top level (outside case branches):
    $\mathsf{OSeq}$ vs.\ $\mathsf{OOp}$,
    $\mathsf{OOp}$ vs.\ $\mathsf{OMap}$,
    $\mathsf{OQuotient}$ vs.\ $\mathsf{OOp}$, etc.
\end{itemize}

\begin{theorem}[NEQ Detection Soundness]
\label{thm:neq-soundness}
Each NEQ certificate $P(a, b)$ is a valid certificate of
non-equivalence: $P(a, b) \implies \den{a} \neq \den{b}$.
\end{theorem}

\begin{proof}
Effect-fiber: a program with IO effects returns
$(\mathsf{None}, \mathit{side\text{-}effect})$ while a pure
program returns $(\mathit{value}, \varnothing)$; these differ
on the heap component.
Literal: different constants trivially differ.
Argument-order: non-commutativity provides a witness
(e.g., $3 - 5 \neq 5 - 3$).
Comparison-operator: $\mathsf{eq}$ uses \texttt{\_\_eq\_\_} while
$\mathsf{is}$ uses pointer identity; \texttt{[] == []} is True
but \texttt{[] is []} is False.
Structural: different-length sequences cannot be equal;
different fold operators compute different functions
(e.g., $\sum \neq \prod$).
Cross-type at top level: an $\mathsf{OSeq}$ (tuple/list literal)
and an $\mathsf{OOp}$ (function application) produce different
types.
\end{proof}

An important conservatism: cross-type certificates are
\emph{suppressed inside case branches}
(\texttt{in\_case\_branch = True}).  Inside a conditional branch,
different structures can compute the same value under the guard's
domain constraint.  For example, $Q[\mathsf{perm}](\mathsf{set}(x))$
and $\mathsf{sorted}(x)$ could agree when
$|\mathsf{set}(x)| \leq 1$.


%% ===============================================================
\section{The Intensional--Extensional Gap}
\label{sec:int-ext-gap}
%% ===============================================================

\subsection{Where the Pipeline Fails}

For 43 of the 50 equivalent benchmark pairs, the pipeline returns
$\mathsf{UNKNOWN}$ or a spurious $H^1$ obstruction.  These failures
cluster into three categories:

\begin{enumerate}
  \item \textbf{Genuinely different algorithms} (8 pairs):
    Graham scan vs.\ monotone chain (convex hull),
    Tarjan vs.\ Kosaraju (SCC),
    matrix exponentiation vs.\ fast doubling (Fibonacci),
    multiplicative formula vs.\ Pascal's triangle ($\binom{n}{k}$),
    Wagner--Fischer vs.\ Hirschberg (edit distance),
    recursive descent vs.\ shunting-yard (expression evaluation),
    stack-based vs.\ recursive (bracket matching),
    trial division vs.\ SPF sieve (prime factorization).
    No structural axiom connects them; they live in disconnected
    components of the rewrite groupoid.

  \item \textbf{Compilation lossy-ness} ($\sim$22 pairs):
    the OTerm compiler loses parameter references, produces
    \texttt{OUnknown} nodes, or generates canonical forms too
    short to trust.

  \item \textbf{Timeout / memory} ($\sim$13 pairs):
    the Z3 fiber checker runs out of time or memory before
    resolving all fibers.
\end{enumerate}

Category~1 is the \emph{fundamental} limitation.  Categories~2 and~3
are engineering problems.

\subsection{Formal Gap Statement}

\begin{theorem}[The Gap]
\label{thm:gap}
There exist programs $f, g : A \to B$ such that:
\begin{enumerate}
  \item $\prod_{x:A} \Path_B(f(x), g(x))$ is inhabited
    (extensionally equal),
  \item $\Path_{\mathsf{PyComp}}(f, g)$ is empty in the current
    HIT (no finite chain of D1--D24 constructors connects them).
\end{enumerate}
\end{theorem}

\begin{proof}
By Rice's theorem: semantic equivalence of programs is undecidable.
Any decidable set of syntactic axioms can only capture a proper
subset of all valid equivalences.  Concretely: Graham scan and
monotone chain compute the same convex hull, but no structural
axiom relates polar-angle sorting to lexicographic sorting, nor
a single-pass stack to a two-pass upper/lower construction.
\end{proof}

\subsection{What the System Cannot Do}

We now formalize the limitations precisely.  Define the
\emph{provability horizon}:

\begin{definition}[Provability Horizon]
\label{def:provability-horizon}
The provability horizon of the pipeline is the set:
\[
  \mathcal{H} \coloneqq \bigl\{(f, g) : f \equiv g \text{ and }
  \Pi(f, g) = \mathsf{EQ}\bigr\}
\]
The \emph{gap set} is:
\[
  \mathcal{G} \coloneqq \bigl\{(f, g) : f \equiv g \text{ and }
  \Pi(f, g) \neq \mathsf{EQ}\bigr\}
\]
\end{definition}

The 8 ``genuinely different algorithm'' pairs lie in $\mathcal{G}$
for a precise reason: each pair consists of two algorithms that
compute the same mathematical function via fundamentally different
decompositions of the problem.  No finite combination of the
syntactic axioms D1--D24 connects the two decompositions, because
the axioms operate on local OTerm structure while the equivalence
is a global property of the algorithm's correctness proof.

\begin{proposition}[Gap Characterization]
\label{prop:gap-char}
A pair $(f, g)$ lies in $\mathcal{G}$ iff the OTerm normal forms
$\mathcal{N}(\den{f})$ and $\mathcal{N}(\den{g})$ lie in
different connected components of the rewrite groupoid
$\mathcal{G}_{\mathsf{PyComp}}$ AND the spec identification
function $\iota$ fails on at least one of them.
\end{proposition}


%% ===============================================================
\section{Enriching the Path Space}
\label{sec:enriching-paths}
%% ===============================================================

\subsection{The Rewrite Groupoid}

\begin{definition}[Rewrite Groupoid]
\label{def:rewrite-groupoid}
$\mathcal{G}_{\mathsf{PyComp}}$ has:
\begin{itemize}
  \item \textbf{Objects}: OTerms in normal form (after $\mathcal{N}$).
  \item \textbf{Morphisms}: finite composites of path constructors
    and their inverses:
    \[
      \mathrm{Hom}_{\mathcal{G}}(a, b) =
      \left\{ p_1^{\epsilon_1} \cdot p_2^{\epsilon_2} \cdots
      p_k^{\epsilon_k} :
      p_i \in \mathrm{Axioms},\; \epsilon_i \in \{+1, -1\}
      \right\} / {\sim}
    \]
    where $\sim$ is the groupoid relations ($p \cdot p^{-1} =
    \mathrm{id}$, associativity).
\end{itemize}
\end{definition}

\begin{proposition}[Fundamental Obstruction]
\label{prop:fundamental-obstruction}
Two OTerms are in different connected components of
$\mathcal{G}_{\mathsf{PyComp}}$ iff no finite chain of axioms
in the equational theory $E$ relates them.  The $\pi_0$ invariant
$\pi_0(\mathcal{G}_{\mathsf{PyComp}})$ classifies these
obstructions.
\end{proposition}

The strategy for improving the prover is to \emph{reduce}
$|\pi_0(\mathcal{G})|$ by adding new axioms --- new path
constructors that connect previously disconnected components.

\subsection{Three Mechanisms for New Paths}

\paragraph{Mechanism 1: Higher-level normalization.}
Every new normalization rule is a new path constructor.
If we add $t \leadsto t'$ to the normalizer, we implicitly add
$\mathsf{norm} : \Path(t, t')$ to the HIT.
Currently missing normalizations include:
\begin{itemize}
  \item \textbf{Accumulator extraction}: recognizing that a
    tail-recursive function with an accumulator computes the same
    value as a non-tail version.
  \item \textbf{Loop interchange}: two nested folds over independent
    ranges can be reordered ($\pi_0$ merge for DP table fills).
  \item \textbf{Divide-and-conquer normalization}:
    $f(a \cdot b) = f(a) \oplus f(b)$ for homomorphisms.
\end{itemize}

\paragraph{Mechanism 2: Specification identification (D20 Yoneda).}
Each new spec pattern $\pi \mapsto S$ creates a new family of paths:
\[
  \forall a, b.\;\; \iota(a) = (S, \vec{x}) \;\wedge\;
  \iota(b) = (S, \vec{x})
  \;\implies\; \Path_{\mathsf{PyComp}}(a, b)
\]
Adding recognition patterns for GCD, LCM, convex hull, edit
distance, topological sort, and shortest path would directly
connect the 8 ``genuinely different algorithm'' pairs.

\paragraph{Mechanism 3: Transport along type coercion paths.}
Univalence provides paths between equivalent types.  If two
programs operate on isomorphic types, transport along the type
equivalence path can normalize them to a common representation.


%% ===============================================================
\section{Strengthening the D20 Spec Engine}
\label{sec:strengthening-d20}
%% ===============================================================

\subsection{The Spec Library as a Presheaf}

\begin{definition}[Spec Presheaf]
$\mathcal{S} : \mathcal{C}^{\mathrm{op}} \to \mathsf{Set}$ assigns
to each type fiber $\tau$ the set of specifications recognizable
on $\tau$.  Restriction maps narrow specifications to sub-fibers.
\end{definition}

The sheaf condition on $\mathcal{S}$ tells us when local spec
matches on individual fibers are sufficient for a global spec
match.

\subsection{Observable Properties as Spec Discovery}

Beyond direct pattern matching, specifications can be
\emph{inferred} from observable properties:

\begin{enumerate}
  \item \textbf{Output type agreement}:
    $\forall \vec{p}.\; \mathsf{tag}(t_f) = \mathsf{tag}(t_g)$.
  \item \textbf{Base-case agreement}: for each $v \in \{0, 1, -1\}$,
    $\mathsf{simplify}(t_f[\vec{p} \mapsto v])
    = \mathsf{simplify}(t_g[\vec{p} \mapsto v])$.
  \item \textbf{Recurrence structure}: alpha-equivalent
    $\mathsf{RecFunction}$ bodies.
  \item \textbf{Idempotence}: $t_f(t_f(x)) = t_f(x)$.
  \item \textbf{Monotonicity}: $x < y \implies t_f(x) \leq t_f(y)$.
\end{enumerate}

\begin{theorem}[Observable Properties as Abstract Spec]
\label{thm:obs-spec}
If two programs agree on output type, all base cases, and their
RecFunction recurrence equations are alpha-equivalent, then they
satisfy the same deterministic specification.  By
Theorem~\ref{thm:det-spec}, they are extensionally equal.
\end{theorem}


%% ===============================================================
\section{Improving the Denotational--Cohomological Interaction}
\label{sec:denot-cohom-interaction}
%% ===============================================================

The current pipeline runs the denotational layer (OTerm) and the
cohomological layer (Z3 + \v{C}ech) \emph{sequentially}.  This
leaves power on the table.

\subsection{Shared Symbol Congruence}

\begin{theorem}[Shared Symbol Congruence]
\label{thm:shared-equiv}
Let $\sigma$ be a shared function symbol.  If
$t_1 |_\tau = t_2 |_\tau$ for every type fiber $\tau$, then
$\sigma(t_1) = \sigma(t_2)$.
\end{theorem}

This principle enables \emph{descent}: decompose a global equality
into local equalities on sub-expressions.

\subsection{Feeding Denotational Evidence into Fibers}

When the denotational layer discovers partial information,
it can be \emph{injected} as lemmas into the Z3 context.

\begin{definition}[Spec-Constrained Fiber]
Given a Z3 term $t$ over parameters $\vec{p}$, a fiber
assignment $\tau$, and discovered properties $P$, the
\emph{spec-constrained fiber check} queries:
\[
  \mathsf{solver.check}\!\left(
    t_f |_\tau \neq t_g |_\tau \;\wedge\;
    \bigwedge_{p \in P} \llbracket p \rrbracket
  \right)
  \stackrel{?}{=} \mathsf{unsat}
\]
\end{definition}


%% ===============================================================
\section{The Graded Verdict}
\label{sec:fallback-graded}
%% ===============================================================

\begin{definition}[Graded Verdict]
\label{def:graded-verdict}
A graded verdict is a tuple $(v, c)$ where
$v \in \{\mathsf{EQ}, \mathsf{NEQ}, \mathsf{UNKNOWN}\}$
and $c = h_0 / \max(|\mathsf{Fibers}|, 1) \in [0,1]$.
Full confidence ($c = 1$) is a global section ---
a proof of equivalence.
\end{definition}

\subsection{H$^1$ Suppression and the Uninterpreted-Depth Filter}

When Z3 compilation produces uninterpreted functions, the fiber
checker may find spurious $H^1$ obstructions.  The implementation
handles this via \texttt{\_uninterp\_depth}: Z3 structural
equality is only trusted when all terms are \emph{fully interpreted}
(depth 0).  Additionally, SAT witnesses with high uninterpreted
depth ($d_1 + d_2 \geq 4$, $\min(d_1, d_2) \geq 1$) are treated
as inconclusive rather than genuine counterexamples.

\begin{definition}[Spurious H$^1$]
An H$^1$ obstruction is \emph{spurious} if it arises from a fiber
where the Z3 terms for $f$ and $g$ differ only in uninterpreted
function symbols.
\end{definition}

The principled response is \emph{refinement}: each builtin
transitioning from uninterpreted to interpreted eliminates a
class of spurious obstructions.


%% ===============================================================
\section{Soundness and Completeness}
\label{sec:fallback-soundness}
%% ===============================================================

\begin{theorem}[Pipeline Soundness]
\label{thm:pipeline-soundness}
The full pipeline $\Pi$ is \emph{sound}: it never returns
$\mathsf{EQ}$ for genuinely non-equivalent programs or
$\mathsf{NEQ}$ for genuinely equivalent programs.
\end{theorem}

\begin{proof}
Each layer $\pi_k$ returns a positive verdict only via a sound
proof:
\begin{itemize}
  \item $\pi_0$ (closed-term evaluation): the denotation of a
    zero-argument function is unique; equal values imply equal
    programs.  Variadic-only calls are trusted only for NEQ
    (witness-based).
  \item $\pi_1$(a) (exact match): identical canonical forms after
    sound normalization rules (Proposition~\ref{prop:phase-soundness})
    imply identical semantics.  The five gating conditions suppress
    potential true positives but never admit false positives.
  \item $\pi_1$(b) (flexible match): OFix keys are hash-based
    identifiers of loop bodies; differing keys do not affect the
    computation when the structural match succeeds on all other
    nodes.
  \item $\pi_1$(c) (HIT + BFS): each path step is a valid axiom
    (Theorem~\ref{thm:axiom-soundness}); composition of valid
    axioms preserves equivalence.
  \item $\pi_1$(d) (provable NEQ): each certificate is sound
    (Theorem~\ref{thm:neq-soundness}).
  \item $\pi_2$ (Z3 structural): trusted only when
    $\texttt{\_uninterp\_depth} = 0$ for all terms.
  \item $\pi_3$ (Z3 per-fiber): UNSAT on a fiber $\tau$ is a proof
    that $f|_\tau = g|_\tau$.  By the stalk criterion
    (Theorem~\ref{thm:stalk-criterion}), agreement on all fibers
    implies global equality.  NEQ requires a concrete SAT witness
    or structural impossibility.
\end{itemize}
\end{proof}

\begin{theorem}[Inherent Incompleteness]
\label{thm:incompleteness}
The pipeline $\Pi$ is \emph{incomplete}: there exist equivalent
programs for which $\Pi(f,g) = \mathsf{UNKNOWN}$.  This is
unavoidable by Rice's theorem.
\end{theorem}

The graded verdict provides a principled measure of ``how
incomplete'' the answer is: the confidence $c$ quantifies
the fraction of the type fibration verified.


%% ===============================================================
\section{The Programme for Improvement}
\label{sec:programme}
%% ===============================================================

\subsection{Axis 1: Richer Normalization (More Paths)}

Each new normalization rule is a new path constructor.
Key additions:
\begin{enumerate}
  \item \textbf{Compiler completeness}: handle all Python
    constructs appearing in the benchmark (Category~2 failures).
  \item \textbf{Accumulator normalization}:
    $\mathsf{fix}[f](\mathsf{acc}, n) \leadsto
    \fold(f_{\mathsf{step}}, \mathsf{acc}_0, \mathsf{range})$.
  \item \textbf{Guard subsumption strengthening}: if
    $G_1 \implies G_2$, simplify
    $\mathsf{case}(G_1, A, \mathsf{case}(G_2, B, C))$ to
    $\mathsf{case}(G_1, A, C)$ in the $G_1$-false branch.
\end{enumerate}

\subsection{Axis 2: Richer Specifications (More Yoneda)}

The path forward is \emph{compositional} spec recognition
mirroring the operadic structure of Chapter~\ref{ch:operads}:
\begin{itemize}
  \item Fold over range with multiplication $\to$
    product-of-range
  \item Product-of-range from 1 to $n$ $\to$ factorial
  \item Product-of-range from $(n-k+1)$ to $n$, divided by
    factorial$(k)$ $\to$ binomial
\end{itemize}
The recognizer becomes a homomorphism from the OTerm operad to
the spec operad.

\subsection{Axis 3: Z3 Compilation Fidelity (Fewer Spurious $H^1$)}

Key builtins to formalize as Z3 RecFunctions:
\begin{itemize}
  \item \texttt{sorted}: axiomatize as a permutation that is
    non-decreasing.
  \item \texttt{set operations}: axiomatize in terms of membership.
  \item \texttt{string operations}: slicing, joining, splitting.
\end{itemize}

\subsection{Concrete Improvement Theorems}

We state specific theorems that, if proved and implemented,
would close identified gaps:

\begin{conjecture}[Accumulator Extraction]
\label{conj:accumulator}
For any tail-recursive function $f$ with linear accumulator
$\alpha$, there exists a fold $\fold(\mathit{step}, \alpha_0, R)$
such that $f \equiv \fold(\mathit{step}, \alpha_0, R)$.
Implementing this as a normalizer rule would connect
$\sim$5 additional benchmark pairs.
\end{conjecture}

\begin{conjecture}[Compositional Spec Recognition]
\label{conj:comp-spec}
The spec identification function $\iota$ can be extended to a
homomorphism $\iota^* : \mathcal{O}_{\mathsf{operad}} \to
\mathcal{S}_{\mathsf{operad}}$ such that $\iota^*$ recognizes
composite specifications (e.g., binomial = ratio of factorials).
This would add $\sim$3 spec families (GCD, LCM, binomial via
multiplicative formula).
\end{conjecture}

\begin{conjecture}[Z3 RecFunction Completeness]
\label{conj:z3-complete}
Axiomatizing \texttt{sorted}, \texttt{set}, and \texttt{str}
operations as Z3 RecFunctions would reduce
$\texttt{\_uninterp\_depth}$ to 0 for $\sim$15 additional
benchmark pairs, converting Z3 timeouts to definitive verdicts.
\end{conjecture}


%% ===============================================================
\section{Summary: What the Theory Tells Us to Build Next}
\label{sec:fallback-summary}
%% ===============================================================

\begin{center}
\begin{tabular}{lll}
\hline
\textbf{Theory} & \textbf{Implementation Target} & \textbf{Expected Impact} \\
\hline
New path constructors & 7-phase normalizer rules & Category 2 failures \\
Yoneda (D20) & Spec library patterns & Category 1 failures \\
Transport & Type coercion paths & Cross-type pairs \\
Sheaf descent & Fiber-aware spec matching & Combined evidence \\
Galois connection & Spec-constrained Z3 & Z3 timeout pairs \\
RecFunction analysis & Z3 recurrence matching & Recursive pairs \\
Operadic composition & Compositional spec recognition & Complex specs \\
\hline
\end{tabular}
\end{center}

Each row is a \emph{theorem} in CCTT that motivates a specific
engineering task.  The theory does not just justify the current
implementation --- it tells us precisely where to push next and
why each extension is sound.
